% =========================================================================
% Capítulo 2: Cinética de la partícula
% Tema 01: Dinámica de la partícula
% Problema 23: Bloque con impulso y pista circular
% Origen: Ayudantía 2, 2026
% =========================================================================
\subsubsection{Bloque con impulso y pista circular}

Un bloque de masa $m$ está inicialmente en reposo sobre una superficie horizontal sin fricción cuando es golpeado por un impulso horizontal $\hat{P}$, como se muestra en la Figura \ref{fig:c2_e01_p23_bloque_impulso}. Después de aplicar el impulso, el bloque se desliza a lo largo de la superficie hasta llegar a una pista circular sin fricción de radio $r$. Sabiendo que la gravedad actúa hacia abajo, determine la velocidad del bloque inmediatamente después de la aplicación del impulso $\hat{P}$.

\begin{figure}[!ht]
    \centering
    \begin{tikzpicture}[>=Latex, scale=1.0]
        % --- PARÁMETROS ---
        \def\r{1.5}      % radio de la pista [u]
        \def\xc{4.0}     % posición del centro de la pista [u]

        % --- SUELO ---
        \fill[pattern=north east lines] (-0.5,-0.25) rectangle (7.5,0);
        \draw[thick] (-0.5,0) -- (7.5,0);

        % --- PISTA CIRCULAR ---
        \draw[thick] (\xc,\r) circle (\r);
        \filldraw (\xc,\r) circle (1.2pt);
        \draw[->] (\xc,\r) -- ({\xc+\r*cos(35)},{\r+\r*sin(35)}) node[midway, above left] {$r$};
        \node[font=\small] at (1.8,2.9) {pista circular};
        \draw (2.5,2.75) -- (3.2,2.4);

        % --- BLOQUE E IMPULSO ---
        \draw[thick, fill=gray!20] (0.9,0) rectangle (1.5,0.6);
        \node at (1.2,0.3) {$m$};
        \draw[->, very thick, myred] (0.0,0.3) -- (0.85,0.3) node[midway, above] {$\hat{P}$};

        % --- GRAVEDAD ---
        \draw[->, very thick] (6.8,2.8) -- (6.8,2.0) node[midway, right] {$\vec{g}$};
    \end{tikzpicture}
    \caption{Bloque golpeado por un impulso, frente a una pista circular.}
    \label{fig:c2_e01_p23_bloque_impulso}
\end{figure}

% --- SOLUCIÓN ---------------------------------------------------------------
\begin{solucion}{Bloque con impulso y pista circular}

\textbf{Velocidad inmediatamente después del impulso}\\
Aplicando el teorema del impulso y la cantidad de movimiento lineal sobre el bloque en reposo:
$$\int \sum \vec{F}\, dt = \Delta \vec{p}$$
$$\hat{P} = m v_1 - m(0) \implies v_1 = \frac{\hat{P}}{m}$$
\respuesta{$\displaystyle v_1 = \frac{\hat{P}}{m}$}

\patoojo{La pista circular y la gravedad no intervienen en esta parte: el impulso es horizontal, el bloque está sobre una superficie sin fricción y las fuerzas finitas (peso, normal) no aportan impulso durante el golpe.}

\end{solucion}
